% Tiny contract
% Teams: General Electronics
% Requirements: REQ-DES-020, REQ-DES-030, REQ-DES-050
% Status: in-progress


The rover's electronics architecture is divided into 5 functional
subsystems: \textit{navigation}, \textit{suspension}, and \textit{arm},
\textit{ground station}, \textit{scientific
payload module},
all coordinated by a central compute unit --- the NVIDIA Jetson Orin 16GB.
The Jetson aggregates sensor data from the navigation subsystem (ToF cameras,
stereo camera, LiDAR, and human presence sensors) via USB and receives
high-level commands from an onboard Raspberry~Pi~4 responsible for arm
control. Low-level motor control for the suspension wheels is delegated to
an ESP32-S3 microcontroller communicating over a CAN bus. (It will also be indicated for the science module and the ground station module. It will also be indicated for the red button and activity indicator, power will be added to the circuit)


\begin{figure}[H]
  \centering
  \adjustbox{max width=\linewidth}{%
    \input{name}%
  }
  \caption{
System Architecture: Electronics Block Diagram}
  \label{fig:block_diagram}
\end{figure}

\begin{table}[H]
  \centering
  \caption{Component model reference
           (see Table~\ref{tab:component_specs} for full specifications)}
  \label{tab:component_models}
  \small
  \begin{tabular}{ll}
    \toprule
    \textbf{Component (diagram label)} & \textbf{Model} \\
    \midrule
    ToF Camera 1--3       & Sipeed MaixSense A010 \\
    Stereo Camera         & StereoLabs ZED2i \\
    LiDAR                 & RPLIDAR 2s \\
    Human Sensor 1--4     & HLK-LD2450 \\
    Camera (arm)          & Arducam AR0234 (2.3\,MP, Global Shutter) \\
    Motor 1--3 (arm)      & SteadyWin GIM8115-36 \\
    Motor 4--6 (arm)      & DAMIAO DM-J4340-2EC \\
    Wheel 1--4            & DAMIAO DM-J10010L-2EC \\
    Wheel 5--8            & SteadyWin GIM8115-36 \\
    NEMA 17 Stepper       & NEMA 17 + driver (TBD) \\
    \bottomrule
  \end{tabular}
\end{table}
